\chapter{Taurus Molecular Cloud}
The Taurus molecular cloud is a nearby star-forming region, located at a distance of approximately 140 parsecs from Earth. It is one of the most studied molecular clouds in the Milky Way, and is known for its rich population of young stars and protostars. The Taurus molecular cloud is composed primarily of molecular hydrogen gas, along with dust and other molecules. It has a complex structure, with filaments, clumps, and cores that are sites of ongoing star formation. The cloud is also home to a variety of astronomical phenomena, including outflows, jets, and Herbig-Haro objects. Studying the Taurus molecular cloud provides valuable insights into the processes of star formation and the early stages of stellar evolution.